% ============================================================
% SECTION 3 — ORIGINAL ARCHITECTURE (THE STARTING POINT)
% ============================================================

\section{Original Architecture}

\begin{tcolorbox}[summarybox={\iconGear[1.5] The Baseline System}]
\color{textdark}
Before the multi-batch hardening process began, RecoverAI V2 existed as a functional but \textbf{fragile} monolithic application. While the core ML and LLM decision-making loops were operational, the surrounding financial infrastructure lacked the rigorous safety boundaries required for a production environment. 

This section documents the \textit{original} architectural state of the system, highlighting the design choices that were later identified as security or reliability risks.
\end{tcolorbox}

\vspace{0.8cm}

% --- VISUALIZING THE ORIGINAL MONOLITH ---
\subsection{System Diagram (Original State)}

The initial architecture was highly coupled. The API layer, orchestration logic, and gateway interactions were not strictly bounded.

\begin{center}
\begin{tikzpicture}[
    node distance=0.8cm and 1.5cm,
    api/.style={rectangle, draw=secondary, thick, fill=bglight, minimum width=3.5cm, minimum height=1cm, rounded corners=4pt, font=\small\bfseries, text=textdark, drop shadow=black!5},
    logic/.style={rectangle, draw=primary, thick, fill=bglight, minimum width=3.5cm, minimum height=1cm, rounded corners=4pt, font=\small\bfseries, text=textdark, drop shadow=black!5},
    db/.style={cylinder, draw=accent, thick, fill=bglight, shape border rotate=90, aspect=0.25, minimum width=3cm, minimum height=1.5cm, font=\small\bfseries, text=textdark, drop shadow=black!5},
    gateway/.style={rectangle, draw=textdark, thick, fill=textdark!10, minimum width=3cm, minimum height=1cm, rounded corners=4pt, font=\small\bfseries, text=textdark, dashed},
    arrow/.style={-{Stealth[scale=1.1]}, thick, draw=textmuted},
    dangerarrow/.style={-{Stealth[scale=1.1]}, thick, draw=danger, dashed}
]

    % Top layer: API
    \node[api] (router) {Monolithic Router (\filepath{router.py})};
    
    % Middle layer: Logic
    \node[logic, below left=1.2cm and -1cm of router] (orchestrator) {Recovery Orchestrator};
    \node[logic, below right=1.2cm and -1cm of router] (refunds) {Refund Endpoint Logic};
    
    % Background worker
    \node[api, left=1cm of router] (reconcile) {Asyncio While-Loop\\(\filepath{main.py})};
    
    % Data layer
    \node[db, below=3cm of router] (sqlite) {SQLite DB};
    
    % External Gateway
    \node[gateway, right=1.5cm of orchestrator, yshift=-1cm] (mock) {Razorpay Mock Service};

    % Connections
    \draw[arrow] (router) -- node[left, font=\scriptsize] {BackgroundTasks} (orchestrator);
    \draw[arrow] (router) -- node[right, font=\scriptsize] {Direct Call} (refunds);
    \draw[arrow] (reconcile) |- (sqlite);
    \draw[arrow] (orchestrator) -- (sqlite);
    \draw[arrow] (refunds) -- (sqlite);
    
    % Danger connections (Direct coupling)
    \draw[dangerarrow] (orchestrator) -- node[above, font=\scriptsize, text=danger] {Unprotected Call} (mock);
    \draw[dangerarrow] (refunds) -- node[right, font=\scriptsize, text=danger] {Unprotected Call} (mock);
    \draw[dangerarrow] (mock) -- (sqlite);
    
\end{tikzpicture}
\end{center}

\begin{tcolorbox}[dangerbox={Critical Architectural Flaws Displayed Above}]
\begin{enumerate}
    \item \textbf{No Execution Guard:} The Orchestrator and the Refund logic directly instantiated and called the Gateway Service without passing through a centralized safety boundary.
    \item \textbf{In-Process Concurrency:} The Reconciliation worker ran as an \code{asyncio} while-loop inside the main thread, risking duplicate sweeps if multiple API pods were booted.
    \item \textbf{Monolithic Routing:} All endpoints (payments, refunds, webhooks) were crammed into a single \filepath{router.py}, making middleware enforcement difficult.
\end{enumerate}
\end{tcolorbox}

\vspace{0.8cm}

% --- COMPONENT BREAKDOWN ---
\subsection{Original Component Breakdown}

\subsubsection{1. The API Layer (\filepath{app/api/router.py})}
Initially, the entire API was defined in a single file. There was no modularity.
\begin{itemize}
    \item \textbf{Authentication:} Non-existent for some endpoints (like WebSockets) and hardcoded to a static string (\code{test\_secret\_key\_123}) for others.
    \item \textbf{Rate Limiting:} None. The system was vulnerable to basic DoS or aggressive API polling.
    \item \textbf{CORS:} Set to \code{allow\_origins=["*"]}, permitting any domain to interface with the backend.
    \item \textbf{Exception Handling:} Broad \code{try...except} blocks swallowed critical database errors (such as \code{IntegrityError} on duplicate transaction inserts).
\end{itemize}

\vspace{0.4cm}

\subsubsection{2. The Orchestrator (\filepath{app/services/orchestrator.py})}
The core brain of the recovery system was highly coupled to the API execution context.
\begin{itemize}
    \item \textbf{Background Execution:} Triggered via FastAPI's \code{BackgroundTasks}. If the Uvicorn/Gunicorn process died during orchestration, the task was permanently dropped from memory with no persistence.
    \item \textbf{Gateway Coupling:} It directly instantiated \code{RazorpayMockService}. There was no abstract \code{GatewayInterface}, making it impossible to easily swap payment providers or safely mock tests without monkeypatching.
    \item \textbf{Execution Logic:} It executed the financial recovery action and \textit{then} attempted to update the database. If a crash occurred in the millisecond between gateway success and database commit, the system would blindly retry the payment upon restart.
\end{itemize}

\vspace{0.4cm}

\subsubsection{3. Refunds Handling}
Refunds were a complete afterthought in the original architecture.
\begin{itemize}
    \item \textbf{No State Machine:} Refunds were triggered via the router, which naively set \code{refund\_status = "REFUND\_PROCESSING"} and fired a background task that blindly updated it to \code{REFUNDED} after a sleep.
    \item \textbf{No Concurrency Protection:} Two rapid requests to refund the same transaction would both read the state as "success" and issue parallel gateway calls.
\end{itemize}

\vspace{0.8cm}

% --- THE DATABASE STATE ---
\subsection{The Database State (Original)}

The original system used \textbf{SQLite} exclusively. While adequate for early prototyping, SQLite fundamentally lacks support for true row-level locking via \code{SELECT ... FOR UPDATE}.

\begin{tcolorbox}[codebox={Original (Flawed) Financial Lock Implementation}]
\begin{lstlisting}[language=Python]
# Original orchestrator logic (Pseudo-code)
txn = db.query(Transaction).filter(Transaction.id == tx_id).first()

if txn.recovery_status == "NOT_STARTED":
    # ⚠️ RACE CONDITION: Two threads can reach here simultaneously
    txn.recovery_status = "EXECUTING"
    db.commit()
    
    # ⚠️ NO EXECUTION GUARD: Direct call
    result = RazorpayMockService.charge(txn)
    
    txn.recovery_status = result.status
    db.commit() 
    # ⚠️ CRASH WINDOW: If process dies before this commit, 
    # the transaction remains stuck or gets double-charged.
\end{lstlisting}
\end{tcolorbox}

Because SQLite silently ignores \code{with\_for\_update()}, the developers were under the false impression that their database queries were locked and safe from concurrent races. It wasn't until the Final Adversarial Audit (post-Batch 4) that the system's reliance on OCC (Optimistic Concurrency Control) for attempts, but lack of OCC for refunds, was fully exposed.

\vspace{0.8cm}

% --- LACK OF IDEMPOTENCY ---
\subsection{The Idempotency Void}

A defining feature of the original architecture was the lack of strict \textbf{system-wide idempotency}.

\begin{center}
\begin{tikzpicture}[
    node distance=1.5cm,
    client/.style={circle, draw=primary, thick, fill=bglight, minimum size=1.5cm, align=center, font=\small\bfseries},
    api/.style={rectangle, draw=secondary, thick, fill=bglight, minimum width=3cm, minimum height=1cm, rounded corners=4pt},
    gateway/.style={rectangle, draw=accent, thick, fill=bglight, minimum width=3cm, minimum height=1cm, rounded corners=4pt},
    arrow/.style={-{Stealth[scale=1.1]}, thick, draw=textdark},
    cross/.style={draw=danger, ultra thick, -}
]

    \node[client] (client) {Client};
    \node[api, right=of client] (api) {RecoverAI API};
    \node[gateway, right=of api] (gateway) {Payment Gateway};
    
    \draw[arrow] (client.45) -- node[above, font=\scriptsize] {Request A} (api.145);
    \draw[arrow] (client.-45) -- node[below, font=\scriptsize] {Request A (Retry)} (api.215);
    
    \draw[arrow] (api.35) -- node[above, font=\scriptsize] {Charge 1} (gateway.145);
    \draw[arrow] (api.-35) -- node[below, font=\scriptsize] {Charge 2 (Duplicate)} (gateway.215);
    
    \node[text=danger, font=\bfseries] at (3.5, 0) {\huge $\times$};
    \node[text=danger, font=\scriptsize, align=center, yshift=-1.2cm] at (3.5, 0) {No Idempotency-Key\\Header Enforcement};

\end{tikzpicture}
\end{center}

Clients could send identical requests, and because the API didn't strictly hash and lock based on an \code{Idempotency-Key} header mapped against a dedicated \code{IdempotencyRecord} table, network retries could theoretically result in duplicate recoveries.

\vspace{0.8cm}

% --- TRANSITION TO THE BATCHES ---
\subsection{The Mandate for Hardening}

By the time the core ML and LLM features were proven viable, it became clear that the underlying engineering infrastructure could not support a financial production environment. 

To resolve this, the engineering team embarked on a rigorous, phased hardening process. The rules were strict: \textbf{No new ML/LLM features could be added until the financial safety, concurrency, and reliability invariants were mathematically proven.}

This hardening effort was divided into \textbf{Batches}, which are detailed in the next section.
